% Part: first-order-logic
% Chapter: models-theories
% Section: expressing-props-of-structures

\documentclass[../../../include/open-logic-section]{subfiles}

\begin{document}

\olfileid{fol}{mat}{exs}

\olsection{Eigenschappen van \printtoken{P}{structure} uitdrukken}

\begin{explain}
Het is vaak nuttig en belangrijk om voorwaarden voor functies en
relaties te formuleren, of algemener, om uit te drukken dat de functies
en relaties in een structuur aan zulke voorwaarden voldoen. We willen
bijvoorbeeld de !!{structure}s voor een taal die ``vatten'' wat de
!!{predicate}s volgens ons moeten ``betekenen'' kunnen onderscheiden
van de !!{structure}s waarvoor dat niet geldt. Uiteraard staat het ons
geheel vrij te bepalen welke !!{structure}s we ``beogen.'' We kunnen
bijvoorbeeld eisen dat de interpretatie van het !!{predicate}~$\le$
een ordening is, of ons beperken tot interpretaties van~$\Lang L$
waarvan het domein uit verzamelingen bestaat en waarin $\Obj \in$ wordt
geïnterpreteerd als de relatie ``is !!a{element} van.'' Maar kunnen we
dit doen met !!{sentence}s van de taal? Anders gezegd: welke
voorwaarden aan !!a{structure}~$\Struct M$ kunnen we uitdrukken door
!!a{sentence} (of misschien een verzameling !!{sentence}s) in de taal
van~$\Struct M$? Sommige voorwaarden zullen we niet kunnen uitdrukken.
Er bestaat bijvoorbeeld geen zin van~$\Lang L_A$ die in
!!a{structure}~$\Struct M$ alleen waar is als $\Domain M = \Nat$. We
kunnen niet uitdrukken dat ``het domein uitsluitend natuurlijke
getallen bevat.'' Maar misschien kunnen we wel ``structurele
eigenschappen'' van !!{structure}s uitdrukken. Welke eigenschappen van
!!{structure}s kunnen we door !!{sentence}s uitdrukken? Of, anders
geformuleerd: welke verzamelingen !!{structure}s kunnen we beschrijven
als de !!{structure}s die !!a{sentence} (of een verzameling
!!{sentence}s) waar maken?
\end{explain}

\begin{defn}[Model van een verzameling]
Laat $\Gamma$ een verzameling !!{sentence}s in een taal~$\Lang L$ zijn.
We zeggen dat !!a{structure}~$\Struct M$ \emph{een model is van}~$\Gamma$
als $\Sat{M}{!A}$ voor alle $!A \in \Gamma$.
\end{defn}

\begin{ex}
De zin $\lforall[x][x \le x]$ is waar in~$\Struct M$ dan en slechts
dan als $\Assign{\le}{M}$ een reflexieve relatie is. De zin
$\lforall[x][\lforall[y][((x \le y \land y \le x) \lif x = y)]]$ is
waar in~$\Struct M$ dan en slechts dan als $\Assign{\le}{M}$
antisymmetrisch is. De zin
$\lforall[x][\lforall[y][\lforall[z][((x \le y \land y \le z)
      \lif x \le z)]]]$ is waar in~$\Struct M$ dan en slechts dan als
$\Assign{\le}{M}$ transitief is. De modellen van
\begin{align*}
\{\quad &\lforall[x][x \le x], \\
   & \lforall[x][\lforall[y][((x \le y \land y \le
    x) \lif x = y)]], \\
   &\lforall[x][\lforall[y][\lforall[z][((x \le y
      \land y \le z) \lif x \le z)]]] \quad \}
\end{align*}
zijn dus precies de structuren waarin~$\Assign{\le}{M}$ reflexief,
antisymmetrisch en transitief is, oftewel een partiële ordening. We
kunnen deze zinnen daarom als axioma's nemen voor de
\emph{eerste-ordetheorie van partiële ordeningen}.
\end{ex}

\end{document}
